\documentclass[10pt,a4paper]{article}

\usepackage[T1]{fontenc}
\usepackage[utf8]{inputenc}
\usepackage[french]{babel}
\usepackage{geometry}
\geometry{top=0.8cm,bottom=2.5cm,left=2.5cm,right=2.5cm}
\usepackage{lmodern}
\usepackage{microtype}
\usepackage{parskip}
\setlength{\parskip}{6pt}
\usepackage{graphicx}
\usepackage{tabularx}
\usepackage{array}
\usepackage{booktabs}
\usepackage{fancyhdr}
\usepackage{hyperref}
\hypersetup{colorlinks=false,pdfborder={0 0 0}}
\usepackage{enumitem}
\usepackage{titlesec}

% ---------- section style ----------
\titleformat{\section}
  {\normalsize\bfseries}
  {}{0em}{}[\vspace{-2pt}\rule{\linewidth}{0.8pt}\vspace{2pt}]
\titlespacing{\section}{0pt}{12pt}{6pt}
\titleformat{\subsection}
  {\normalsize\bfseries\itshape}{}{0em}{}
\titlespacing{\subsection}{0pt}{8pt}{3pt}

% ---------- footer ----------
\pagestyle{fancy}
\fancyhf{}
\renewcommand{\headrulewidth}{0pt}
\fancyfoot[C]{\tiny
  Generali Iard, Société anonyme au capital de 94~630~300 euros ---
  Entreprise régie par le Code des assurances --- 552~062~663 RCS Paris
  --- IDU ADEME FR232327\_01NBYI\\
  Generali Vie, Société anonyme au capital de 341~059~488 euros ---
  Fonds de Retraite Professionnelle Supplémentaire régi par le Code des assurances
  --- 880~265~418 RCS Paris --- IDU ADEME FR232327\_01NBYI\\
  Siège social : 89 rue Taitbout --- 75009 Paris --- Sociétés appartenant
  au Groupe Generali immatriculé sur le registre italien des groupes
  d'assurances sous le numéro 026}

\begin{document}
\thispagestyle{fancy}

% ================================================================
%  TITLE
% ================================================================
\begin{center}
  {\Large\bfseries Présentation de l'entreprise et du service de R\&D}\\[6pt]
  {\normalsize\itshape Dossier de candidature CIFRE --- Yassine MANNAI}
\end{center}

\vspace{0.4cm}

% ================================================================
\section{Présentation de Generali France}
% ================================================================

Generali France est l'un des principaux assureurs de l'Hexagone, s'appuyant
sur l'esprit pionnier d'un groupe international fondé en 1831.  L'entreprise
opère sur l'ensemble des segments de l'assurance (Vie, Dommages, Santé,
Prévoyance et Retraite) avec une ambition de \og{}Lifetime Partner\fg{} pour
ses clients.

% ----------------------------------------------------------------
\subsection{Indicateurs clés et volume d'activité (2025)}
% ----------------------------------------------------------------
% NOTE : chiffres FY2025 — à confirmer sur le rapport annuel Generali France publié début 2026.
\begin{itemize}[leftmargin=1.6em,itemsep=1pt]
  \item \textbf{Chiffre d'affaires :} 20,5 milliards d'euros.                   % TODO: confirmer
  \item \textbf{Résultat opérationnel :} 1~400 millions d'euros (normes IFRS~17/9). % TODO: confirmer
  \item \textbf{Actifs sous gestion :} 124 milliards d'euros.                   % TODO: confirmer
  \item \textbf{Effectif total :} 9~300 collaborateurs et salariés commerciaux.
  \item \textbf{Rayonnement :} 8 millions de clients et un réseau de près de 800 agents généraux.
\end{itemize}

% ----------------------------------------------------------------
\subsection{Organisation et direction de l'encadrement}
% ----------------------------------------------------------------

L'encadrement industriel de la thèse est rattaché à la \textbf{Direction des
Investissements} de Generali France, placée sous la responsabilité de
\textbf{Cédrik d'Aviau de Ternay} (Directeur des Investissements).  Cette
direction pilote l'allocation stratégique du bilan de Generali France dans le
cadre de Solvabilité~II, en s'appuyant notamment sur l'outil propriétaire du
Groupe, \textbf{GPS} (\textit{Group Portfolio Solutions}).

% ================================================================
\section{Équipe d'accueil : Allocation Stratégique d'Actifs (SAA)}
% ================================================================

Le candidat sera intégré au sein de l'\textbf{équipe Allocation Stratégique
d'Actifs (SAA)}, entité de recherche quantitative et de pilotage du bilan au
sein de la Direction des Investissements.  L'équipe regroupe des expertises en
finance quantitative, gestion actif-passif (ALM), modélisation actuarielle et
optimisation de portefeuille sous contraintes réglementaires.

\begin{itemize}[leftmargin=1.6em,itemsep=1pt]
  \item \textbf{Directeur des Investissements :} Cédrik d'Aviau de Ternay.
  \item \textbf{Responsable scientifique (Encadrant industriel) :}
        Mamy Ramamonjisoa, membre de l'équipe SAA.
  \item \textbf{Doctorant CIFRE (candidat) :} Yassine Mannai.
\end{itemize}

% ================================================================
\section{Le service de R\&D : équipe Allocation Stratégique d'Actifs}
% ================================================================

L'équipe SAA constitue le pôle de recherche quantitative appliquée de la
Direction des Investissements.  Elle est chargée de :

\begin{itemize}[leftmargin=1.6em,itemsep=1pt]
  \item \textbf{Pilotage stratégique du bilan :} construction et révision de
        l'allocation stratégique d'actifs dans le cadre de Solvabilité~II, via
        l'outil propriétaire GPS (\textit{Group Portfolio Solutions}).
  \item \textbf{Modélisation des risques :} développement et maintenance du
        générateur de scénarios économiques (ESG), calcul du capital réglementaire
        (SCR) par simulation imbriquée (LSMC), et optimisation actif-passif (ALM).
  \item \textbf{Recherche et innovation :} intégration des approches de la
        finance quantitative contemporaine (modèles génératifs, apprentissage par
        renforcement, proxy ML) dans les processus d'allocation et de couverture.
  \item \textbf{Coordination opérationnelle :} transmission des instructions de
        réallocation à Generali Asset Management (GenAM) et suivi de l'exécution.
\end{itemize}

% ================================================================
\section{Projet de thèse CIFRE}
% ================================================================

% ----------------------------------------------------------------
\subsection{Titre}
% ----------------------------------------------------------------

\begin{center}
  \itshape
  Dépendance factorielle conditionnelle aux régimes, proxy de capital supervisé
  et couverture profonde pour le Total Portfolio Approach en assurance
\end{center}

% ----------------------------------------------------------------
\subsection{Contexte et motivation industrielle}
% ----------------------------------------------------------------

La gestion du risque de dépendance entre facteurs économiques dans un contexte
de changements de régime constitue un enjeu croissant pour les assureurs.
L'épisode de 2022 --- au cours duquel la corrélation entre actions et
obligations s'est inversée simultanément sur tous les marchés --- a révélé les
limites des approches classiques d'allocation par classes d'actifs et motive le
développement d'une approche \textbf{factorielle, dynamique et intégrée} au
sein de l'équipe SAA.

Les outils actuels de l'équipe souffrent de quatre limitations structurelles :
(i) des corrélations calibrées sur des moyennes historiques, aveugles aux
ruptures de régime ; (ii) un proxy LSMC polynomial précis uniquement dans la
région de calibration, dégradé précisément dans les zones de stress les plus
critiques ; (iii) une décision d'allocation et de couverture effectuée en
silos, sans politique jointe ; (iv) des actifs privés à valorisation lissée
introduisant un biais de sous-estimation du risque réel.

La thèse vise à fournir les fondations quantitatives permettant de rendre
opérationnel le \textbf{Total Portfolio Approach (TPA)} pour un assureur sous
Solvabilité~II.

% ----------------------------------------------------------------
\subsection{Axes de recherche}
% ----------------------------------------------------------------

\begin{enumerate}[leftmargin=2em,itemsep=4pt]
  \item \textbf{Simulation génératrice conditionnelle aux régimes (Axe~1).}
        Développement et comparaison de modèles génératifs
        (GAN conditionnel, VAE conditionnel, Score-Based Generative Models)
        pour simuler la distribution jointe des facteurs TPA sous chaque régime
        macroéconomique identifié en temps réel par un filtre séquentiel.
        Validation formelle de la structure de dépendance apprise par des tests
        de spécification sur copules conditionnelles (Fermanian, 2013, 2022)
        et établissement de garanties de convergence Wasserstein-2.

  \item \textbf{Proxy de SCR supervisé pour la décision dynamique (Axe~2).}
        Construction d'un proxy du capital réglementaire (SCR)
        exprimé en fonction des facteurs TPA, du régime courant et des états
        des passifs (dont les lags de NAV d'actifs privés), offrant un temps
        d'inférence inférieur à 10~ms et des bornes d'erreur contrôlées
        par régime --- condition nécessaire à son intégration comme contrainte
        explicite dans une boucle d'optimisation dynamique.

  \item \textbf{Politique jointe d'allocation et de couverture profonde (Axe~3).}
        Apprentissage d'une politique de décision dynamique optimisant
        conjointement les expositions aux facteurs TPA et les positions de
        couverture (\textit{deep hedging}), avec une fonction objectif
        identifiée automatiquement à partir du facteur de risque dominant,
        sous contraintes de capital (proxy SCR), de liquidité et d'engagements
        en fonds fermés.
\end{enumerate}

% ----------------------------------------------------------------
\subsection{Encadrement}
% ----------------------------------------------------------------

\begin{tabularx}{\textwidth}{@{}lX@{}}
  \toprule
  \textbf{Rôle}                  & \textbf{Personne / Institution} \\
  \midrule
  Candidat                       & Yassine Mannai \\
  Directrice de thèse            & Caroline Hillairet
                                   (Laboratoire CREST, ENSAE Paris) \\
  Co-directeur de thèse          & Anthony Réveillac
                                   (INSA Toulouse,
                                    Institut de Mathématiques de Toulouse) \\
  Encadrant industriel           & Mamy Ramamonjisoa ---
                                   équipe SAA, Generali France \\
  Directeur des Investissements  & Cédrik d'Aviau de Ternay ---
                                   Generali France \\
  Durée                          & 36 mois (CDD) \\
  Démarrage prévu                & Octobre -- Janvier 2026/2027 \\
  \bottomrule
\end{tabularx}

% ================================================================
\section{Intérêt des travaux pour le développement de l'entreprise}
% ================================================================

La thèse s'inscrit dans le pilier \og{}Excellence\fg{} du plan stratégique
\textbf{Boost 2027}.  Pour Generali France, l'intérêt est triple :

\begin{enumerate}[leftmargin=2em,itemsep=4pt]
  \item \textbf{Résilience de l'allocation stratégique.}
        Les travaux fourniront un outil de simulation capable d'anticiper les
        ruptures de régime et d'adapter l'allocation du bilan avant les phases
        de stress, réduisant l'exposition aux pertes extrêmes dues à des
        dépendances non linéaires entre facteurs.

  \item \textbf{Accélération du calcul du capital réglementaire.}
        Le proxy SCR supervisé permettra d'intégrer le coût en capital
        directement dans les boucles d'optimisation en temps quasi-réel,
        rendant possible une SAA dynamique --- là où le recalcul complet
        (Prophet/MoSes) nécessite aujourd'hui plusieurs heures.

  \item \textbf{Opérationnalisation du Total Portfolio Approach.}
        Les modules développés seront conçus pour être intégrés à \textbf{GPS}
        (\textit{Group Portfolio Solutions}), l'outil propriétaire du Groupe,
        offrant à Generali France un avantage concurrentiel durable dans la
        gestion actif-passif et le pilotage du bilan sous Solvabilité~II.
\end{enumerate}

\end{document}
